\documentclass[12pt,a4paper]{article}
\usepackage[utf8]{inputenc} % inutile avec XeLaTeX/LuaLaTeX
\usepackage[T1]{fontenc}
\usepackage{amsmath,amssymb,mathrsfs,tikz,times,pifont}
\usepackage{enumitem}
\usepackage{multicol}
\usepackage{lmodern}
\newcommand\circitem[1]{%
\tikz[baseline=(char.base)]{
\node[circle,draw=gray, fill=red!55,
minimum size=1.2em,inner sep=0] (char) {#1};}}
\newcommand\boxitem[1]{%
\tikz[baseline=(char.base)]{
\node[fill=cyan,
minimum size=1.2em,inner sep=0] (char) {#1};}}
\setlist[enumerate,1]{label=\protect\circitem{\arabic*}}
\setlist[enumerate,2]{label=\protect\boxitem{\alph*}}
\everymath{\displaystyle}
\usepackage[left=1cm,right=1cm,top=1cm,bottom=1.7cm]{geometry}
\usepackage[colorlinks=true, linkcolor=blue, urlcolor=blue, citecolor=blue]{hyperref}
\usepackage{array,multirow}
\usepackage[most]{tcolorbox}
\usepackage{varwidth}
\usepackage{float}
\tcbuselibrary{skins,hooks}
\usetikzlibrary{patterns}

\usepackage{pgfplots}
\pgfplotsset{compat=1.15}
\usepackage{mathrsfs}

\usepackage{amsmath}
\usepackage{amsfonts}
\usepackage{amssymb}
\usepackage{tkz-tab}
\usepackage{mdframed}
\usepackage{tikz}
\usetikzlibrary{arrows, shapes.geometric, fit}

\usepackage{xcolor}
\usepackage{tikz}
\usepackage{tkz-tab}
\usepackage{verbatim} % INDISPENSABLE pour l'environnement comment
\usepackage{yhmath}
\usepackage{array,multirow}
\usepackage[most]{tcolorbox}
\usepackage{varwidth}
\usepackage{graphicx}
\usetikzlibrary{shapes.geometric, positioning}
\usepackage{tabularx}
\usepackage{pifont} % Pour les symboles de coches et croix
\usepackage{eso-pic}         % Pour ajouter des éléments en arrière-plan
\usetikzlibrary{angles,quotes,arrows.meta}
% Définition du rouge pour les titres et soulignements
\definecolor{monRouge}{RGB}{180, 0, 0}

\usepackage{graphicx}
\usetikzlibrary{trees} % Nécessaire pour l'arbre
\usetikzlibrary{shapes.geometric, positioning}
\usepackage{tabularx}
\usepackage{pifont} % Pour les symboles de coches et croix
\usepackage{enumitem}
\usepackage{makecell}
% Définition des symboles du tableau
\newcommand{\coche}{{\color{red}\ding{51}}}
\newcommand{\croix}{{\color{red}\ding{55}}}

% Commande pour ajouter du texte en arrière-plan
\AddToShipoutPicture{
    \AtTextCenter{%
        \makebox[0pt]{\rotatebox{45}{\textcolor[gray]{0.6}{\fontsize{5cm}{5cm}\selectfont PGB}}}
    }
}
\renewcommand{\seriesdefault}{b} % rend tout plus foncé

\begin{document}

\tiny

\textbf{Exercice D'application} \\[2pt]

Soit l'urne contenant: 4 boules noires, 3 boules vertes et 2 boules rouges.


On tire 4 boules de l'urne. Déterminer le nombre de cas où l'on obtient :

\begin{enumerate}
    \item Exactement une verte parmi les boules tirées.
    \item Au moins deux vertes parmi les boules tirées.
    \item Au plus une rouge parmi les boules tirées.
    \item Aucune boule noire.
    \item Une rouge puis deux noires puis une verte.
\end{enumerate}

\textbf{Solution} \\[2pt]

\renewcommand{\arraystretch}{2} % Pour aérer un peu les lignes du tableau

\raggedright
\begin{tabular}{|p{6cm}|p{6cm}|p{5cm}|}
\hline
\textbf{Simultané} & \textbf{Successif sans remise} & \textbf{Successif avec remise}\\
\hline
\makecell[l]{\textbf{1)} il faut tirer : $V \ \overline{V} \ \overline{V} \ \overline{V}$\\ le nombre demandé est :$C_3^1 \times C_6^3 $} & \makecell[l]{\textbf{1)} il faut tirer :$ V \ \overline{V} \ \overline{V} \ \overline{V}$\\ en tenant l'ordre en compte\\ le nombre demandé est :\\$ A_3^1 \times A_6^3 \times \color{red} C_4^1 \times C_3^3$} & \makecell[l]{\textbf{1)} il faut tirer :$ V \ \overline{V} \ \overline{V} \ \overline{V}$\\ en tenant l'ordre en compte\\ le nombre demandé est :\\$ 3^1 \times 6^3 \times \color{red} C_4^1 \times C_3^3 $} \\
\hline
\makecell[l]{\textbf{2)} il faut tirer :\\ $V V \overline{V} \overline{V} \text{ ou } V V V \overline{V}$\\ le nombre demandé est :\\$ C_3^2 \times C_6^2 + C_3^3 C_6^1 $} & \makecell[l]{\textbf{2)} il faut tirer : $V V \overline{V} \ \overline{V} \text{ ou } V V V \overline{V}$\\ en tenant compte de l'ordre.\\ le nombre demandé est :\\$ A_3^2 \times A_6^2 \times \color{red} C_4^2 \times C_2^2 \color{black} +$ \\$ A_3^3 \times A_6^1 \times \color{red} C_4^3 \times C_1^1 $} & \makecell[l]{\textbf{2)} il faut tirer :\\$ V V V V \text{ ou } V V V \overline{V} \text{ ou }$\\$ V V \overline{V} \ \overline{V} $\\ en tenant l'ordre en compte\\ le nombre demandé est :\\$ 3^4 + 3^3 \times 6^1 \times \color{red} C_4^3 \times C_1^1 \color{black} + $\\$3^2 \times 6^2 \times \color{red} C_4^2 \times C_2^2$} \\
\hline
\makecell[l]{\textbf{3)} il faut tirer :\\ $R \overline{R} \ \overline{R} \ \overline{R} \text{ ou } \overline{R} \ \overline{R} \ \overline{R} \ \overline{R}$\\ le nombre demandé est :\\$ C_2^1 \times C_7^3 + C_7^4 $} & \makecell[l]{\textbf{3)} il faut tirer :\\ $R \overline{R} \ \overline{R} \ \overline{R} \text{ ou } \overline{R} \overline{R} \ \overline{R} \ \overline{R}$\\ en tenant compte de l'ordre.\\ le nombre demandé est :\\$ A_2^1 \times A_7^3 \times \color{red} C_4^1 \times C_3^3 \color{black} + A_7^4$} & \makecell[l]{\textbf{2)} il faut tirer :\\$ R \overline{R} \ \overline{R} \overline{R} \text{ ou } \overline{R} \ \overline{R} \ \overline{R} \ \overline{R} $\\ en tenant l'ordre en compte\\ le nombre demandé est :\\$ 2^1 \times 7^3 \times \color{red} C_4^1 \times C_3^3 \color{black} + 7^4$} \\
\hline
\makecell[l]{\textbf{4)} il faut tirer : $ \overline{N} \ \overline{N} \ \overline{N} \ \overline{N} $\\ le nombre demandé est :$ C_5^4 $} & \makecell[l]{\textbf{4)} il faut tirer : $ \overline{ N} \ \overline{ N} \overline{ N} \ \overline{ N} $\\ le nombre demandé est : $ A_5^4 $} & \makecell[l]{\textbf{4)} il faut tirer : $ \overline{N} \ \overline{N} \ \overline{N} \ \overline{N} $\\ le nombre demandé est : $ 5^4 $} \\
\hline
\makecell[l]{\textbf{5)} \color{red}{IMPOSSIBLE} } & \makecell[l]{\textbf{5)} il faut tirer : $ \overline{R} \ \overline{N} \overline{N} \ \overline{V} $\\ le nombre demandé est :\\ $ A_2^1 \times A_4^2 \times A_3^1 $} & \makecell[l]{\textbf{5)} il faut tirer : $ \overline{R} \ \overline{N} \ \overline{N} \ \overline{V} $\\ le nombre demandé est :\\ $ 2^1 \times 4^2 \times 3^1 $} \\
\hline
\end{tabular}

\vspace{1cm}
--------------------------------------------------------------------------------
\vspace{1cm}

\textbf{Exercice D'application} \\[2pt]

Soit l'urne contenant: 4 boules noires, 3 boules vertes et 2 boules rouges.


On tire 4 boules de l'urne. Déterminer le nombre de cas où l'on obtient :

\begin{enumerate}
    \item Exactement une verte parmi les boules tirées.
    \item Au moins deux vertes parmi les boules tirées.
    \item Au plus une rouge parmi les boules tirées.
    \item Aucune boule noire.
    \item Une rouge puis deux noires puis une verte.
\end{enumerate}

\textbf{Solution} \\[2pt]

\renewcommand{\arraystretch}{2} % Pour aérer un peu les lignes du tableau

\raggedright
\begin{tabular}{|p{6cm}|p{6cm}|p{5cm}|}
\hline
\textbf{Simultané} & \textbf{Successif sans remise} & \textbf{Successif avec remise}\\
\hline
\makecell[l]{\textbf{1)} il faut tirer : $V \ \overline{V} \ \overline{V} \ \overline{V}$\\ le nombre demandé est :$C_3^1 \times C_6^3 $} & \makecell[l]{\textbf{1)} il faut tirer :$ V \ \overline{V} \ \overline{V} \ \overline{V}$\\ en tenant l'ordre en compte\\ le nombre demandé est :\\$ A_3^1 \times A_6^3 \times \color{red} C_4^1 \times C_3^3$} & \makecell[l]{\textbf{1)} il faut tirer :$ V \ \overline{V} \ \overline{V} \ \overline{V}$\\ en tenant l'ordre en compte\\ le nombre demandé est :\\$ 3^1 \times 6^3 \times \color{red} C_4^1 \times C_3^3 $} \\
\hline
\makecell[l]{\textbf{2)} il faut tirer :\\ $V V \overline{V} \overline{V} \text{ ou } V V V \overline{V}$\\ le nombre demandé est :\\$ C_3^2 \times C_6^2 + C_3^3 C_6^1 $} & \makecell[l]{\textbf{2)} il faut tirer : $V V \overline{V} \ \overline{V} \text{ ou } V V V \overline{V}$\\ en tenant compte de l'ordre.\\ le nombre demandé est :\\$ A_3^2 \times A_6^2 \times \color{red} C_4^2 \times C_2^2 \color{black} +$ \\$ A_3^3 \times A_6^1 \times \color{red} C_4^3 \times C_1^1 $} & \makecell[l]{\textbf{2)} il faut tirer :\\$ V V V V \text{ ou } V V V \overline{V} \text{ ou }$\\$ V V \overline{V} \ \overline{V} $\\ en tenant l'ordre en compte\\ le nombre demandé est :\\$ 3^4 + 3^3 \times 6^1 \times \color{red} C_4^3 \times C_1^1 \color{black} + $\\$3^2 \times 6^2 \times \color{red} C_4^2 \times C_2^2$} \\
\hline
\makecell[l]{\textbf{3)} il faut tirer :\\ $R \overline{R} \ \overline{R} \ \overline{R} \text{ ou } \overline{R} \ \overline{R} \ \overline{R} \ \overline{R}$\\ le nombre demandé est :\\$ C_2^1 \times C_7^3 + C_7^4 $} & \makecell[l]{\textbf{3)} il faut tirer :\\ $R \overline{R} \ \overline{R} \ \overline{R} \text{ ou } \overline{R} \overline{R} \ \overline{R} \ \overline{R}$\\ en tenant compte de l'ordre.\\ le nombre demandé est :\\$ A_2^1 \times A_7^3 \times \color{red} C_4^1 \times C_3^3 \color{black} + A_7^4$} & \makecell[l]{\textbf{2)} il faut tirer :\\$ R \overline{R} \ \overline{R} \overline{R} \text{ ou } \overline{R} \ \overline{R} \ \overline{R} \ \overline{R} $\\ en tenant l'ordre en compte\\ le nombre demandé est :\\$ 2^1 \times 7^3 \times \color{red} C_4^1 \times C_3^3 \color{black} + 7^4$} \\
\hline
\makecell[l]{\textbf{4)} il faut tirer : $ \overline{N} \ \overline{N} \ \overline{N} \ \overline{N} $\\ le nombre demandé est :$ C_5^4 $} & \makecell[l]{\textbf{4)} il faut tirer : $ \overline{ N} \ \overline{ N} \overline{ N} \ \overline{ N} $\\ le nombre demandé est : $ A_5^4 $} & \makecell[l]{\textbf{4)} il faut tirer : $ \overline{N} \ \overline{N} \ \overline{N} \ \overline{N} $\\ le nombre demandé est : $ 5^4 $} \\
\hline
\makecell[l]{\textbf{5)} \color{red}{IMPOSSIBLE} } & \makecell[l]{\textbf{5)} il faut tirer : $ \overline{R} \ \overline{N} \overline{N} \ \overline{V} $\\ le nombre demandé est :\\ $ A_2^1 \times A_4^2 \times A_3^1 $} & \makecell[l]{\textbf{5)} il faut tirer : $ \overline{R} \ \overline{N} \ \overline{N} \ \overline{V} $\\ le nombre demandé est :\\ $ 2^1 \times 4^2 \times 3^1 $} \\
\hline
\end{tabular}
\end{document}
